\documentclass[12pt,letterpaper]{article}



%%
%% Required packages:
%%
\usepackage{etex}
\usepackage{amsmath}
\usepackage{amssymb}
\usepackage{amstext}
\usepackage{amsthm}
\usepackage{fancyhdr,fancybox}
\usepackage{mathrsfs} % need for \mathscr command
\usepackage{multicol}
\usepackage[normalem]{ulem}
%\usepackage{pictex,pictexplus}
% \usepackage{pictexwd}
\usepackage{rotating}
\usepackage{url}
\usepackage{xfrac}
\usepackage{booktabs}
\usepackage{color,soul}
\usepackage{comment}

\usepackage{hyperref}
\hypersetup{
    colorlinks=true,     % false: boxed links; true: colored links
    linkcolor=blue,      % color of internal links
    citecolor=blue,      % color of links to bibliography
    filecolor=blue,      % color of file links
    urlcolor=blue        % color of external links
}


\usepackage{tikz}
\usetikzlibrary{backgrounds,calc}

%%
%% Common formatting commands
%%
\setlength{\parindent}{0in}
\setlength{\textwidth}{7in}
\setlength{\evensidemargin}{-0.25in}
\setlength{\oddsidemargin}{-0.25in}
\setlength{\parskip}{.5\baselineskip}
\setlength{\topmargin}{-0.5in}
\setlength{\textheight}{9in}

%               Problem and Part
\newcounter{problemnumber}
\newcounter{partnumber}
\newcounter{subpartnumber}
\newcommand{\Problem}{%
  \refstepcounter{problemnumber}%
  \setcounter{partnumber}{0}%
  \setcounter{subpartnumber}{0}%
  \item[\makebox{\hfil\textbf{\theproblemnumber.}\hfil}]%
  }
\newcommand{\InProblem}[2][1.6]{%
  \refstepcounter{problemnumber}%
  \setcounter{partnumber}{0}%
  \setcounter{subpartnumber}{0}%
  \textbf{\theproblemnumber.}\ \ \parbox[t]{#1in}{#2}%
}
\newcommand{\InSmallProblem}[1]{\InProblem[1.05]{#1}}
\newcommand{\NoProblem}[1]{%
  \mbox{\hfil\textbf{#1}\hfil}%
  }
\newcommand{\UnProblem}{%
  \refstepcounter{problemnumber}%
  \setcounter{partnumber}{0}%
  \setcounter{subpartnumber}{0}%
  \mbox{\hfil\textbf{\theproblemnumber.}\hfil}%
  }
\newcommand{\InFlexProblem}[1]{%
  \refstepcounter{problemnumber}%
  \setcounter{partnumber}{0}%
  \setcounter{subpartnumber}{0}%
  \textbf{\theproblemnumber.}\ \ #1%
}
%%  Part:
\newcommand{\Part}{%
  \renewcommand{\thepartnumber}{(\alph{partnumber})}%
  \refstepcounter{partnumber}
  \setcounter{subpartnumber}{0}%
  \item[(\alph{partnumber})]%
}
\newcommand{\InPart}[2][1.75]{%
  \renewcommand{\thepartnumber}{(\alph{partnumber})}%
  \refstepcounter{partnumber}%
  \setcounter{subpartnumber}{0}%
  (\alph{partnumber})\ \ \parbox[t]{#1in}{#2}%
}
\newcommand{\UnPart}{%
  \refstepcounter{partnumber}%
  \renewcommand{\thepartnumber}{(\alph{partnumber})}%
  \setcounter{subpartnumber}{0}%
  (\alph{partnumber})%
}
\newcommand{\InLargePart}[1]{\InPart[3.05]{#1}}
\newcommand{\InFullPart}[1]{\InPart[6.2]{#1}}
\newcommand{\InSmallPart}[1]{\InPart[1.05]{#1}}
%% SubPart:
\newcommand{\SubPart}{%
  \refstepcounter{subpartnumber}%
  \item[(\roman{subpartnumber})]%
  }
\newcommand{\InSubPart}[2][2.25]{%
  \stepcounter{subpartnumber}%
  (\roman{subpartnumber})\ \ \parbox[t]{#1in}{#2}%
}
\newcommand{\InSmallSubPart}[1]{\InSubPart[1.5]{#1}}
\newcommand{\InTinySubPart}[1]{%
  \stepcounter{subpartnumber}%
  (\roman{subpartnumber})\ \ \makebox{#1}%
}
\newcommand{\InMySubPart}[2][2.25]{%
  \stepcounter{subpartnumber}%
  \parbox[t]{#1in}{\makebox[0.3in][l]{(\roman{subpartnumber})}\ \ {#2}}%
}
\newcommand{\TrueFalse}{\textbf{T}\ \ \textbf{F}\ \ }
\newcommand{\TFPart}[1]{% 
  \stepcounter{partnumber}%
  \setcounter{subpartnumber}{0}%
  \item[(\alph{partnumber})] \TrueFalse\parbox[t]{5.5in}{#1}}

\newcommand{\Points}[1]{(#1 points) \ }
\newcommand{\Point}[1]{(#1 point) \ }


%% For Answers:
\newcounter{answernumber}
\newcommand{\Answer}{%
  \refstepcounter{answernumber}%
  \setcounter{partnumber}{0}%
  \setcounter{subpartnumber}{0}%
  \item[\makebox{\hfil\textbf{\theanswernumber.}\hfil}]%
  }


% Setting up the headers for the first page
%\chead{} % will default to what is typed in the file.
\fancyhead[L]{\textsc{Math 22a}}
\fancyhead[R]{\textsc{Fall 2024}}
\fancyfoot[R]{
	\vspace*{\fill}\footnotesize\textcopyright{} \emph{Harvard Math 22a}	
}
\renewcommand{\headrulewidth}{0pt}

%\fancyhead[C]{}
%	\chead{}	
%	\lhead{}
%	\rhead{}
%	\cfoot{}


% Putting copyright on the footer of each page
%\fancypagestyle{secondpage}{
%	\fancyhead[C]{TESTing again} % change this to be blank
%		%\chead{}	
%	\fancyhead[L]{bears} % change this to be blank
%	\fancyhead[R]{dogs} % change this to be blank
%	\fancyfoot[R]{ %keeps the footer the same
%		\vspace*{\fill}\footnotesize\textcopyright{} \emph{Harvard Math 22a}	
%	}
%}
%\renewcommand{\headrulewidth}{0pt}
%\pagestyle{fancy}
\pagestyle{empty}


%\lhead{}
%\rhead{}
%\rfoot{\vspace*{\fill}\footnotesize\textcopyright{} \emph{Harvard Math 22a}}
%\renewcommand{\headrulewidth}{0pt}
%\pagestyle{fancy}




%%
%% Common shortcut commands:
%%
\newcommand{\ds}{\displaystyle}
\newcommand{\sgn}{\operatorname{sgn}}
\renewcommand{\Pr}{\operatorname{P}}
\newcommand{\CPr}[3][{}]{\Pr_{\text{#1}}\left(#2\,|\,#3\right)}

\newcommand{\C}{\mathbb{C}}
\newcommand{\R}{\mathbb{R}}
\newcommand{\Z}{\mathbb{Z}}
\newcommand{\N}{\mathbb{N}}
\newcommand{\Q}{\mathbb{Q}}
\newcommand{\D}{\mathbb{D}}

\newcommand{\rref}{\operatorname{rref}}
\newcommand{\im}{\operatorname{im}}
\newcommand{\rank}{\operatorname{rank}}
\newcommand{\diag}{\operatorname{diag}}
\newcommand{\Span}{\operatorname{span}}
%\renewcommand{\vec}[1]{\mathbf{#1}}
\newcommand{\charpoly}{\mathscr{P}}
\newcommand{\nicefrac}[2]{\sfrac{#1}{#2}} % using xfrac package
\newcommand{\topology}{\mathcal{T}}
\newcommand{\Inn}{\operatorname{Inn}}

\newcommand{\blankbox}[2]{\fbox{\rule{#1}{0in}\rule{0in}{#2}}}

\newenvironment{myindentpar}[1]%
 {\begin{list}{}%
         {\setlength{\leftmargin}{#1}
          \setlength{\rightmargin}{\leftmargin}}%
         \item[]%
 }
 {\end{list}}
\newtheorem{theorem}{Theorem}
\newtheorem*{NStheorem}{Theorem}
\newtheorem{cor}[theorem]{Corollary}
\newtheorem*{claim}{Claim}
\newtheorem{defn}{Definition}

\newenvironment{amatrix}[1]{%
  \left(\begin{array}{@{}*{#1}{c}|c@{}}
}{%
  \end{array}\right)
}

%--------grstep
% For denoting a Gauss' reduction step.
% Use as: \grstep{\rho_1+\rho_3} or \grstep[2\rho_5 \\ 3\rho_6]{\rho_1+\rho_3}
\newcommand{\grstep}[2][\relax]{%
   \ensuremath{\mathrel{
       {\mathop{\longrightarrow}\limits^{#2\mathstrut}_{
                                     \begin{subarray}{l} #1 \end{subarray}}}}}}
\newcommand{\swap}{\leftrightarrow}


\usetikzlibrary{arrows}
\usepackage{wrapfig}

\makeatletter
\renewcommand*\env@matrix[1][*\c@MaxMatrixCols c]{%
	\hskip -\arraycolsep
	\let\@ifnextchar\new@ifnextchar
	\array{#1}}
\makeatother

\def\env@matrix{\hskip -\arraycolsep
	\let\@ifnextchar\new@ifnextchar
	\array{*\c@MaxMatrixCols c}}

\chead{\Large\textbf{Null Spaces and Column Spaces, $\vec\S$4.2}}% 

\begin{document}
\thispagestyle{fancy}

Recall from last class:

\begin{itemize}
\item A real {\bf vector space} is a non-empty set $V$ on which two operations, scalar multiplication and addition, are defined subject to the following 10 axioms. We call elements in $V$ vectors. Let $u,v,w \in V$ and $c,d\in \mathbb{R}$.
\begin{enumerate}
\item $u+v \in V$.
\item $u+v=v+u$.
\item $(u+v)+w=u+(v+w)$.
\item there exists a vector $0_V$ in $V$ such that $u+0_V=u$.
\item there exists a vector $(-u)$ in $V$ such that $u+(-u)=0_V$.
\item $cu \in V$.
\item $c(u+v)=cu+cv$.
\item $(c+d)u=cu+du$.
\item $c(du)=(cd)u$.
\item $1\cdot u =u$.
\end{enumerate}
\item A {\bf subsapce} of a vector space $V$ is a subset of $V$ such that
\begin{enumerate}
\item $0_V \in H$.
\item $H$ is closed under vector addition; i.e. if $u, v \in H$, then $u+v \in H$.
\item $H$ is closed under scalar multiplication; i.e. if $c \in \mathbb{R}$ and $u \in H$, then $cu \in H$.
\end{enumerate}
\item {\bf Theorem 4.1} If $v_1, \dots, v_p \in V$, then $\text{Span}\{v_1 ,\dots, v_p\}$ is a subspace of $V$. 
\item We call $\text{Span}\{v_1 ,\dots, v_p\}$ the subspace generated by or spanned by $v_1,\dots, v_p$.
\end{itemize}

\newpage

\begin{defn}
	The {\bf null space} of an $m \times n$ matrix $A$ is $$\text{Null}(A)=N(A)=\{x \in \mathbb{R}^n : A\vec{x}=\vec{0}\}.$$
\end{defn}

{\bf Theorem:} If $A$ is an $m\times n$ matrix, then $N(A)$ is a subspace of $\mathbb{R}^n$.

\begin{enumerate}

\Problem Prove the Theorem.

\vfill

\Problem Let $\vec{b}$ be a non-zero element in $\mathbb{R}^m$. Is $\{ \vec{x} \in \mathbb{R}^n : A\vec{x}=\vec{b}\}$ a subspace of $\mathbb{R}^n$?

\vfill

\Problem Let $A=\begin{bmatrix} 1 & -4 & 0 & 2 &0\\  0 & 0 & 1 & -5 & 0\\ 0 & 0 & 0 & 0 & 2 \end{bmatrix}$. Find $N(A)$.
%\Part Let $A=\begin{bmatrix} 1 & 2 & 4 & 0 &0\\  0 & 1 & 3 & -2 & 0\\ \end{bmatrix}$.

%\vfill

%\Part 

\vfill 

\newpage

\begin{defn}
	The {\bf column space} of an $m\times n$ matrix $A$ is the span of the columns of $A$. We denote the column space of $A$ by $\text{Col}(A)$.
\end{defn}

{\bf Theorem:} If $A$ is an $m\times n$ matrix, then $Col(A)$ is a subspace of $\mathbb{R}^m$.

\Problem Let $A=\begin{bmatrix} 1 & -4 & 0 & 2 &0\\  0 & 0 & 1 & -5 & 0\\ 0 & 0 & 0 & 0 & 2 \end{bmatrix}$. Find $\text{Col}(A)$.

\newpage

\begin{defn}
Let $V$ and $W$ be vector spaces, then $T: V \to W$ is a {\bf linear transformation} if 
\begin{enumerate}
	\item $T(u+v)=T(u)+T(v)$ for all $u, v \in V$ and 
	\item $T(c u)=c T(u)$ for all $c \in \mathbb{R}$ and $u \in V$.
\end{enumerate}
\end{defn}

\begin{defn}
	The {\bf range or image } of $T$ is defined by $$\text{Range}(T)=\text{im}(T)=\{T(x): x \in V\}=\{w\in W: \text{there exists } a \in V \text{ such that } T(a)=w\}.$$ 
	
	The {\bf kernel} of $T$ is defined by $$\text{ker}(T)=\{ x \in V: T(x)=0_W\}.$$
\end{defn}

\Problem Define $T: \mathbb{P}_n \to \mathbb{P}_{n-1}$ by $T(f)=\frac{df}{dt}$.
\Part Is $T$ linear?

\vfill

\Part Find the image of $T$.

\vfill

\Part Find the kernel of $T$.

\vfill

\newpage

\Problem Define $T: \mathbb{P}_2 \to \mathbb{P}_3$ by $T(p)=t p(t)+p'(t)$.

\Part Find the range of $T$.

\vfill

\Part Find the kernel of $T$.

\vfill 

\Part Is $T$ onto? Is $T$ one-to-one?

\vfill

\newpage

\Problem Let $T: \mathbb{P}_2 \to \mathbb{R}^2$ by $T(p)=\begin{bmatrix} p(0)\\ p(1)\\ \end{bmatrix}$.

\Part Is $T$ linear?

\vfill

\Part Find $\text{ker}(T)$ and $\text{Range}(T)$.

\vfill 

\newpage

\Problem Define $T: \mathbb{P}_2 \to M_{2 \times2 }$ by $T(p)=\begin{bmatrix}
p(1)-p(0) & 0\\
0 & p(0)\\
\end{bmatrix}$. Find the kernel and image of $T$.

\end{enumerate}

\begin{comment}
\newpage
\chead{\Large\textbf{Null Spaces/Column Spaces -- Answers / Solutions}}%
\lhead{}
\rhead{}
\thispagestyle{fancy}

\begin{enumerate}
  \Answer 
  \begin{proof} Let $A$ be an $m \times n$ matrix, then $N(A)\subseteq \mathbb{R}^n$. We check the three subspace conditions.
	
	{\bf (a)} Since $A \cdot 0 = 0$, $0 \in N(A)$.
	
	{\bf (b)} Let $x, y \in N(A)$, then $A(x+y)=Ax+Ay=0+0=0$. Thus $(x+y) \in N(A)$.
	
	{\bf (c)} Let $x \in N(A)$ and $c \in \mathbb{R}$, then $A(cx)=c(Ax)=c(0)=0$. Thus $cx \in N(A)$.
	
	Thus $N(A)$ is a subspace of $\mathbb{R}^n$.
\end{proof}

\Answer No, the set is not a subspace. Since $A0 \neq b$, $0 \not \in \{x : Ax=b\}$. Thus the first subspace condition fails.

\Answer Fortunately, $A$ is nearly in reduced row echelon form. By definition $\vec{x}\in N(A)$ if and only if $A\vec{x}=\vec{0}$. Now $A\vec{x}=\vec{0}$ if and only if 
\begin{align*}
x_1 - 4x_2 +2 x_4 &=0\\
x_3 -5x_4 &=0\\
2x_5 &= 0.
\end{align*}
Thus $\vec{x} \in N(A)$ if and only if $$\vec{x}=\begin{bmatrix} 4x_2 -2x_4\\ x_2\\ 5x_4\\ x_4\\ 0\\ \end{bmatrix}= x_2 \begin{bmatrix} 4\\ 1\\ 0\\ 0\\ 0\\ \end{bmatrix} + x_4 \begin{bmatrix} -2\\ 0\\ 5\\ 1\\ 0\\ \end{bmatrix}.$$ Thus $\vec{x} \in N(A)$ if and only if $$\vec{x} \in \text{Span}\left\{ \begin{bmatrix} 4\\ 1\\ 0\\ 0\\ 0\\ \end{bmatrix},\begin{bmatrix} -2\\ 0\\ 5\\ 1\\ 0\\ \end{bmatrix}\right\}.$$
Therefore, $$N(A)=\text{Span}\left\{ \begin{bmatrix} 4\\ 1\\ 0\\ 0\\ 0\\ \end{bmatrix},\begin{bmatrix} -2\\ 0\\ 5\\ 1\\ 0\\ \end{bmatrix}\right\}.$$

\Answer The column space of $A$ is defined to be the span of the columns of $A$. Thus $$\text{Col}(A)=\text{Span}\left\{
\begin{bmatrix} 1\\ 0\\ 0\\
\end{bmatrix},
\begin{bmatrix} -4\\ 0\\ 0\\
\end{bmatrix},
\begin{bmatrix} 0\\ 1\\ 0\\
\end{bmatrix},
\begin{bmatrix} 2\\ -5\\ 0\\
\end{bmatrix},
\begin{bmatrix} 0\\ 0\\ 2\\
\end{bmatrix}
\right\}.$$ Notice that some of the vectors in the span are redundant, so we can express the column space of $A$ using only the first, third, and fifth vectors; namely, $$\text{Col}(A)=\text{Span}\left\{
\begin{bmatrix} 1\\ 0\\ 0\\
\end{bmatrix},
\begin{bmatrix} 0\\ 1\\ 0\\
\end{bmatrix},
\begin{bmatrix} 0\\ 0\\ 2\\
\end{bmatrix}
\right\}.$$
\Answer
\begin{enumerate}
	\item We check the two linearity conditions. Let $c \in \mathbb{R}$ and $f, g \mathbb{P}_n$. Note $$T(f+g)= \frac{d}{dt}\left(f+g\right)=\frac{df}{dt}+\frac{dg}{dt}$$ from the sum rule for differentiation. Thus $T(f+g)=T(f)+T(g)$.
	Now we check $T$ respects scalar multiplication: $$T(cf)=\frac{d}{dt}\left(cf\right)=c \frac{df}{dt}$$ from the constant multiple differentiation rule. Thus $T(cf)=cT(f)$
	
	Since the two conditions hold, $T$ is linear.
	\item Suppose $f \in \mathbb{P}_{n-1}$, then $f(t)=a_0 +a_1 t +\dots + a_{n-1} t^{n-1}$. Antidifferentiating we get $T(a_0 t+\frac{a_1}{2} t^2 +\dots + \frac{a_{n-1}}{n}t^n)=f$, and hence $T$ is surjective. Thus the image of $T$ is $\mathbb{P}_{n-1}$.
	
	We could also see this in the following way. Every element in $\mathbb{P}_n$ can be written as a linear combination of $1, t, t^2, \dots, t^n$, and hence every element in the image of $T$ can be written as a linear combination of $T(1), T(t), T(t^2),\dots, T(t^n)$. Thus the image of $T$ is the span of $T(1), T(t), T(t^2),\dots, T(t^n)$, which gives the image of $T$ is the span of $1, 2t,3t^2, \dots, n t^{n-1}$. Thus the image is all of $\mathbb{P}_{n-1}$.
	\item Let $f \in \mathbb{P}_n$, then $f(t)=a_0 +a_1 t+\dots + a_n t^n$. Now $f \in \text{ker}(T)$ if and only if $T(f)=0$. Thus we get $T(f)=a_1+2a_2 t+\dots + n a_n t^{n-1} =0$. Equating coefficients, we get $a_1=0, 2a_2=0, \dots, na_n =0$, and hence $a_1=a_2=\dots=a_n=0$. Thus $a_0$ is a free variable and all the other coefficients must be zero. Thus $f$ is in the kernel if and only if $f(t)=a_0$. Therefore $$\text{ker}(T)= \{a_0 : a_0 \in \mathbb{R}\}=\text{Span}\{1\}.$$ 
\end{enumerate}

\Answer
\begin{enumerate}
	\item Again note that every element in $\mathbb{P}_2$ can be written as a linear combination of $1, t, t^2$, so the image of $T$ will be given by linear combinations of $T(1), T(t), T(t^2)$. Note $T(1)=t$, $T(t)=t^2+1$, and $T(t^2)=t^3+2t$. Thus $$\text{Range}(T)=\text{Span}\{t, t^2+1, t^3+2t\}.$$
	\item Let $f \in \mathbb{P}_2$, then $f(t)=a_0+a_1 t +a_2 t^2$. Now $f \in \text{ker}(T)$ if and only if $T(f)=0$. Now $T(f)= t(a_0+a_1 t+a_2 t^2)+a_1+2a_2 t =a_1 +t(a_0+2a_2)+a_1t^2+a_2 t^3$. Thus $T(f)=0$ if and only if $a_1=0$, $a_0+2a_2=0$, $a_1=0$, and $a_2=0$. Thus all the coefficients of $f$ must be zero, and hence $\text{ker}(T)=\{0_{\mathbb{P}_2}\}$.
	\item 
	
	{\bf Claim:} $T$ is one-to-one.
	
	Let $f, g \in \mathbb{P}_2$, then $f(t)=a_0+a_1t +a_2 t^2$ and $g(t)=b_0+b_1 t +b_2 t^2$. Suppose $T(f)=T(g)$, then $a_1 +t(a_0+2a_2)+a_1t^2+a_2 t^3=b_1 +t(b_0+2b_2)+b_1t^2+b_2 t^3$. Equating coefficients, we get $a_1=b_1$, $a_0+2a_2=b_0+2b_2$, $a_1=b_1$, and $a_2=b_2$. Thus $f=g$, and hence $T$ is one-to-one.
	
	{\bf Claim:} $T$ is not onto.
	
	We show that $1$ is not in the image of $T$. Suppose $1$ is in the image of $T$, then $1 = a t+ b (t^2+1)+c (t^3+2t)=b+(a+2c)t+bt^2+ct^3$ for some scalars $a,b,c$. Thus equating coefficients we get $1=b$, $a+2c=0$, $b=0$, and $c=0$. But we have a contradiction since $b$ cannot be zero and one! Thus $1$ cannot be in the range of $T$ and hence $T$ is not surjective.
\end{enumerate}
\Answer
\begin{enumerate}
	\item Yes. $T$ is linear, since $$T(cp(t)) =\begin{bmatrix} c p(0)\\ c p(1)\\ \end{bmatrix}= c \begin{bmatrix} p(0)\\ p(1)\\ \end{bmatrix} =c T(p)$$ and $$T(p+q)=\begin{bmatrix} (p+q)(0)\\ (p+q)(1)\\ \end{bmatrix} = \begin{bmatrix} p(0)+q(0)\\ p(1)+ q(1) \end{bmatrix} = \begin{bmatrix} p(0)\\ p(1)\\ \end{bmatrix} + \begin{bmatrix} q(0)\\ q(1)\\ \end{bmatrix}= T(p)+T(q).$$
	\item Let $f \in \mathbb{P}_2$, then $f(t)=a_0 +a_1 t+a_2 t^2$. Thus $T(f)=\begin{bmatrix} a_0\\ a_0+a_1+a_2\\ \end{bmatrix}$. Thus $$\text{Range}(T)= \left\{\begin{bmatrix} a_0\\ a_0+a_1+a_2\\ \end{bmatrix} : a_0, a_1, a_2 \in \mathbb{R} \right\}=\text{Span}\left\{\begin{bmatrix} 1\\ 1\\ \end{bmatrix}, \begin{bmatrix} 0\\ 1\\ \end{bmatrix} \right\}=\mathbb{R}^2.$$
	
	Now $f \in \text{ker}(T)$ if and only if $T(f)=0$. Thus $f \in \text{ker}(T)$ if and only if $a_0=0$ and $a_0 +a_1+a_2=0$. Thus $f(t)=-a_2t+a_2t^2$. Hence $$\text{ker}(T)=\{-a_2 t+a_2t^2 : a_2 \in \mathbb{R} \}=\text{Span}\{t(1-t)\}.$$ Thus the kernel of $T$ consists of polynomials of degree 2 or less that have zeros at $t=0$ and $t=1$.
\end{enumerate}
\Answer Note $T(1)=\begin{bmatrix} 0 & 0\\ 0 & 1\\ \end{bmatrix}$, $T(t)=\begin{bmatrix} 1 & 0\\ 0 & 0\\ \end{bmatrix}$, and $T(t^2)=\begin{bmatrix} 1 & 0\\ 0 & 0\\ \end{bmatrix}$. Thus $$\text{Range}(T)=\text{Span}\left\{\begin{bmatrix} 0 & 0\\ 0 & 1\\ \end{bmatrix},\begin{bmatrix} 1 & 0\\ 0 & 0\\ \end{bmatrix} \right\}.$$ Note this is different from $\text{Span}\left\{\begin{bmatrix} 1 & 0\\ 0 & 1\\ \end{bmatrix} \right\}$ since the entries along the diagonal can be different. 

Now we compute the kernel of $T$. Let $f \in \mathbb{P}_2$, then $f(t)=a_0+a_1t+a_2t^2$. Now $f \in \text{ker}(T)$ if and only if $T(f)=0_{M_{n\times n}}$, and hence we get $\begin{bmatrix} a_1+a_2 & 0\\ 0 & a_0\\ \end{bmatrix}=\begin{bmatrix} 0 & 0\\ 0 & 0\\ \end{bmatrix}$. Thus we get $a_0=0$ and $a_1+a_2=0$. Thus $f(t)=-a_2t +a_2 t^2$. Finally $$\text{ker}(T)=\{ -a_2 t+a_2 t^2: a_2 \in \mathbb{R} \}=\text{Span}\{t(1-t) \}.$$
\end{enumerate} 

\end{comment}
\end{document}


%%% Local ables: 
%%% mode: latex
%%% TeX-master: t
%%% End: 
